\documentclass[es]{P3CTeX} % idioma de la asignatura: es (castellano)
\PECTeXconfig{
  data = {
    student = {José Castillo},
    grade   = {Grado de Ingeniería Informática}
  },
  no-plagi
}

\NewDocumentCommand{\LoadPreamble}{O{PR} m m}{
  \PECTeXconfig{
    data = {
      subj-fullname  = {Programación Orientada a Objetos},
      subj-shortname = {POO},
      subj-code      = {05.555},
      ca-shortname   = {#2},
      ca-fullname    = {#3},
      date           = {\today}
    },
    #1
  }
  \PECTeXcover{outter-backtitle, subtitle, subsubtitle}
}
\AtEndDocument{
  \nocite{*}
  \bibliographystyle{abbrvnat}
  \bibliography{../references}
}

\LoadPreamble{PAC1}{Primera Prueba de Evaluación Continua}
\begin{document}

\section{Introducción}
Esta PAC de ejemplo ha sido generada por el plugin \texttt{pecuoc} para demostrar
el flujo de trabajo con la clase \PECTeX. El contenido está en castellano, mientras
que el código y los identificadores se mantienen en inglés.

\section{Ejercicio 1: comparativa de paradigmas}
A continuación se resumen tres pilares de la programación orientada a objetos
mediante el módulo \texttt{pxTAB}:

\pxTABsetup{style=header, align=l}
\pxTable{Concepto, Definición, Mecanismo}%
  {Encapsulación, Ocultar el estado interno, campos privados}%
  {Herencia, Reutilizar comportamiento, palabra clave extends}%
  {Polimorfismo, Una interfaz y varias formas, sobrescritura de métodos}

\section{Ejercicio 2: diseño de clases con UML}
El módulo \texttt{pxUML} permite declarar clases y dibujarlas:

\pxUMLClass{Vehiculo}{- matricula : String}[+ arrancar() : void]
\pxUMLClass{Coche}{- numPuertas : int}[+ arrancar() : void]

\pxUMLDiagram{
  \drawclass[text width=14em]{Vehiculo}[0,0]
  \drawclass[text width=14em]{Coche}[0,-4]
}

\end{document}
